\documentclass[12pt,a4paper]{article}

\usepackage[utf8]{inputenc}
\usepackage[T1]{fontenc}
\usepackage[a4paper,margin=1in]{geometry}
\usepackage{hyperref}
\usepackage{xcolor}
\usepackage{longtable}
\usepackage{array}
\usepackage{enumitem}
\usepackage{graphicx}
\usepackage{booktabs}
\usepackage{tabularx}
\usepackage{float}
\usepackage{listings}

\hypersetup{
    colorlinks=true,
    linkcolor=blue,
    urlcolor=blue
}

\definecolor{codebg}{RGB}{245,245,245}
\definecolor{codeframe}{RGB}{220,220,220}

\lstset{
    basicstyle=\ttfamily\small,
    backgroundcolor=\color{codebg},
    frame=single,
    rulecolor=\color{codeframe},
    breaklines=true,
    columns=fullflexible,
    keepspaces=true
}

\newcounter{reportfigure}
\newcommand{\reportfigure}[3]{
    \refstepcounter{reportfigure}
    \begin{center}
        \includegraphics[width=#1]{#2}\\[4pt]
        \small\textbf{Figure \thereportfigure.} #3
    \end{center}
}

\title{Current Progress Report\\ESP32 Adaptive Sampling IoT Project}
\author{IoT Adaptive Sampling Project}
\date{April 18, 2026}

\begin{document}

\maketitle

\section{Introduction}

This document is the current detailed progress report for the project. It is
meant to complement the older historical report in
\texttt{PHASE1\_REPORT.tex}, which documents Phases 1 to 7 in depth but does
not reflect the latest hardware validation, saved MQTT evidence, three-signal
results, or the current project status.

The purpose of this report is to present:

\begin{itemize}
    \item what has been implemented so far,
    \item what has already been validated on real hardware,
    \item what measured results exist today,
    \item what is still pending before final submission.
\end{itemize}

\section{Executive Summary}

The project now has a working \texttt{ESP32 + FreeRTOS} pipeline that:

\begin{itemize}
    \item generates the required input signal through a virtual sensor,
    \item collects samples into fixed-duration windows,
    \item estimates the dominant frequency from each window,
    \item adapts the sampling rate at runtime,
    \item computes a per-window aggregate,
    \item publishes real aggregate messages to a nearby edge server over
    \texttt{MQTT/WiFi},
    \item prepares compact aggregate payloads for \texttt{LoRaWAN + TTN}.
\end{itemize}

The strongest completed milestone is the end-to-end \texttt{MQTT/WiFi}
validation on the real Heltec board. Real timing data were captured with
synchronized timestamps, and three different signal profiles were also
validated on the real board.

The main items still pending are:

\begin{itemize}
    \item the live \texttt{LoRaWAN/TTN} campus validation,
    \item energy measurements with a cleaner two-board setup,
    \item one live \texttt{MQTTS} validation run if a TLS-capable broker is available.
\end{itemize}

\section{Current Phase Status}

\begin{longtable}{|>{\raggedright\arraybackslash}p{0.08\textwidth}|>{\raggedright\arraybackslash}p{0.23\textwidth}|>{\raggedright\arraybackslash}p{0.14\textwidth}|>{\raggedright\arraybackslash}p{0.43\textwidth}|}
\hline
\textbf{Phase} & \textbf{Topic} & \textbf{Status} & \textbf{Notes} \\
\hline
\texttt{1} & Environment setup & Complete & Board, PlatformIO, serial, and broker environment validated \\
\hline
\texttt{2} & Firmware skeleton & Complete & Modular FreeRTOS structure and shared data types implemented \\
\hline
\texttt{3} & Virtual sensor & Complete & Required sinusoidal signal generated in firmware \\
\hline
\texttt{4} & Maximum sampling frequency & Complete & Strict benchmark executed on the real board \\
\hline
\texttt{5} & Frequency analysis & Complete & Dominant-frequency detection validated \\
\hline
\texttt{6} & Adaptive sampling & Complete & Runtime update from \texttt{50 Hz} to \texttt{40 Hz} validated \\
\hline
\texttt{7} & Aggregate computation & Complete & Window averages computed and passed downstream \\
\hline
\texttt{8} & MQTT over WiFi & Complete & End-to-end home-network validation completed on hardware \\
\hline
\texttt{9} & LoRaWAN + TTN & Partial & Payload path implemented, but real TTN proof still pending \\
\hline
\texttt{10} & Performance measurements & Partial & Latency and payload-size evidence exist; energy still pending \\
\hline
\texttt{11} & Optional physical sensor / extras & Optional & Not required for the core project \\
\hline
\texttt{12} & Final packaging and write-up & In progress & README, evidence map, and screenshot assets are being prepared \\
\hline
\end{longtable}

\section{Project Architecture}

The implemented data path is:

\begin{lstlisting}
virtual signal
    -> sampling task
    -> FFT / dominant-frequency detection
    -> adaptive sampling control
    -> 5 s aggregate window
    -> MQTT over WiFi to nearby edge listener
    -> compact LoRaWAN payload for TTN
\end{lstlisting}

The modular structure is important because each task has a clear
responsibility:

\begin{itemize}
    \item \texttt{signal\_input}: sample generation and runtime sampling-rate changes,
    \item \texttt{signal\_processing}: window buffering, frequency estimation, and aggregate creation,
    \item \texttt{sampling\_control}: adaptive policy,
    \item \texttt{comm\_mqtt}: edge publishing,
    \item \texttt{comm\_lorawan}: compact uplink payload preparation,
    \item \texttt{metrics}: runtime counters and timing support.
\end{itemize}

\section{Hardware And Software Used}

\subsection{Hardware}

\begin{itemize}
    \item \texttt{Heltec WiFi LoRa 32 V3} compatible board
    \item integrated \texttt{WiFi}
    \item integrated \texttt{868 MHz LoRa} radio
    \item integrated OLED
    \item laptop running the local broker and edge listener
\end{itemize}

\subsection{Software}

\begin{itemize}
    \item \texttt{PlatformIO}
    \item \texttt{ESP-IDF}
    \item \texttt{Mosquitto}
    \item Python-based \texttt{MQTT} listener and analyzer
    \item \texttt{BetterSerialPlotter}
\end{itemize}

\section{Signal Model And Profiles}

The core virtual input signal follows the assignment requirement:

\begin{lstlisting}
2*sin(2*pi*3*t) + 4*sin(2*pi*5*t)
\end{lstlisting}

This produces two spectral components:

\begin{itemize}
    \item \texttt{3 Hz}
    \item \texttt{5 Hz}
\end{itemize}

The current implementation uses three profiles:

\begin{center}
\begin{tabularx}{\textwidth}{|>{\raggedright\arraybackslash}p{0.22\textwidth}|>{\raggedright\arraybackslash}p{0.18\textwidth}|X|}
\hline
\textbf{Profile} & \textbf{Purpose} & \textbf{Notes} \\
\hline
\texttt{clean\_reference} & baseline & deterministic signal without noise \\
\hline
\texttt{noisy\_reference} & robustness & same signal plus low random noise \\
\hline
\texttt{anomaly\_stress} & bonus / stress & noisy signal plus sparse spike injection \\
\hline
\end{tabularx}
\end{center}

The selected profile is carried through the aggregate object and the MQTT
payload as \texttt{signal\_profile}, and anomaly windows also report
\texttt{anomaly\_count}.

\section{Maximum Sampling Frequency Benchmark}

The strict benchmark from Phase 4 tested multiple target rates and classified a
rate as stable only if timing stayed clean under the project rule.

\subsection{Measured Result Table}

\begin{center}
\small
\begin{tabularx}{\textwidth}{|>{\raggedright\arraybackslash}p{0.12\textwidth}|>{\raggedright\arraybackslash}p{0.14\textwidth}|>{\centering\arraybackslash}p{0.08\textwidth}|>{\centering\arraybackslash}p{0.09\textwidth}|>{\centering\arraybackslash}p{0.09\textwidth}|X|}
\hline
\textbf{Target} & \textbf{Achieved} & \textbf{Drops} & \textbf{Miss.} & \textbf{Stable} & \textbf{Observation} \\
\hline
\texttt{50 Hz} & \texttt{50.25 Hz} & 0 & 0 & Yes & clean stable run \\
\hline
\texttt{100 Hz} & \texttt{100.25 Hz} & 0 & 1 & No & very close, but one miss occurred \\
\hline
\texttt{200 Hz} & \texttt{200.25 Hz} & 0 & 2 & No & throughput held, timing did not \\
\hline
\texttt{250 Hz} & \texttt{250.25 Hz} & 0 & 12 & No & watchdog instability started \\
\hline
\texttt{500 Hz} & \texttt{500.25 Hz} & 0 & 26 & No & repeated instability \\
\hline
\texttt{1000 Hz} & \texttt{1000.00 Hz} & 1 & 55 & No & severe instability \\
\hline
\end{tabularx}
\normalsize
\end{center}

\subsection{Interpretation}

The important conclusion is:

\begin{itemize}
    \item the board can technically generate and consume data above \texttt{50 Hz},
    \item but under the current strict stability rule, the highest clean operating point is \texttt{50 Hz}.
\end{itemize}

This is why the project uses:

\begin{itemize}
    \item \texttt{50 Hz} as the baseline oversampling reference,
    \item \texttt{40 Hz} as the first stable adaptive rate for the current signal.
\end{itemize}

\section{FFT And Adaptive Sampling}

After the first stable \texttt{50 Hz} window, the firmware detects the
dominant component correctly and updates the rate:

\begin{lstlisting}
Window 0 ready | fs=50.0Hz samples=250 avg=0.0041 min=-5.6139 max=5.6139 duration=4.97s
Adaptive sampling update | window=0 dominant=5.00 Hz peak=2.0000 previous_fs=50.0 Hz new_fs=40.0 Hz
Applying adaptive sampling frequency 40.0 Hz
\end{lstlisting}

\subsection{Why \texttt{40 Hz} Is Correct}

The current policy uses:

\begin{itemize}
    \item oversampling factor: \texttt{8x}
    \item dominant frequency: \texttt{5 Hz}
\end{itemize}

So:

\begin{lstlisting}
5 Hz * 8 = 40 Hz
\end{lstlisting}

That result matches the validated runtime behavior.

\section{Aggregate Computation}

Each completed window becomes an aggregate object that contains:

\begin{itemize}
    \item \texttt{window\_id}
    \item \texttt{window\_start\_us}
    \item \texttt{window\_end\_us}
    \item \texttt{sample\_count}
    \item \texttt{sampling\_frequency\_hz}
    \item \texttt{dominant\_frequency\_hz}
    \item \texttt{average\_value}
    \item \texttt{signal\_profile}
    \item \texttt{anomaly\_count}
\end{itemize}

The first clean aggregate appeared at the initial baseline rate:

\begin{lstlisting}
Aggregate result | window=0 fs=50.0Hz samples=250 avg=0.0041 dominant=5.00 Hz
\end{lstlisting}

After adaptation, the aggregate pipeline remained correct:

\begin{lstlisting}
Aggregate result | window=1 fs=40.0Hz samples=200 avg=-0.0000 dominant=5.00 Hz
\end{lstlisting}

This confirms that:

\begin{itemize}
    \item the aggregate function is correct,
    \item the adaptive transition does not break the aggregate path,
    \item the edge and cloud paths can consume real summary data instead of placeholders.
\end{itemize}

\section{MQTT Over WiFi Validation}

The local edge path was fully validated on the real board on the home network.

\subsection{What Was Proven}

\begin{itemize}
    \item the Heltec connected to the home WiFi,
    \item the board reached the local broker,
    \item the edge listener received real aggregate messages,
    \item the messages contained synchronized timestamps,
    \item the saved run had no missing windows.
\end{itemize}

\subsection{Latest Clean Listener Run}

Source: \texttt{results/summaries/mqtt\_summary\_2026-04-18\_listener.md}

\subsubsection*{Overview}

\begin{center}
\begin{tabular}{|l|l|}
\hline
\textbf{Metric} & \textbf{Value} \\
\hline
Messages received & \texttt{5} \\
\hline
Unique windows & \texttt{5} \\
\hline
Missing windows & \texttt{0} \\
\hline
First window & \texttt{296} \\
\hline
Last window & \texttt{300} \\
\hline
Received span & \texttt{20.002 s} \\
\hline
Signal profile & \texttt{clean\_reference} \\
\hline
Total payload bytes & \texttt{2270} \\
\hline
\end{tabular}
\end{center}

\subsubsection*{Timing Metrics}

\begin{center}
\begin{tabular}{|l|r|r|r|r|}
\hline
\textbf{Metric} & \textbf{Min} & \textbf{Avg} & \textbf{P95} & \textbf{Max} \\
\hline
\texttt{listener\_latency\_us} & 1232577 & 1234421.6 & 1239997 & 1239997 \\
\hline
\texttt{end\_to\_end\_latency\_us} & 1584866 & 1587868.6 & 1596541 & 1596541 \\
\hline
\texttt{edge\_delay\_us} & 352289 & 353447.0 & 356544 & 356544 \\
\hline
\end{tabular}
\end{center}

\subsubsection*{Signal Metrics}

\begin{center}
\begin{tabular}{|l|r|r|r|}
\hline
\textbf{Metric} & \textbf{Min} & \textbf{Avg} & \textbf{Max} \\
\hline
\texttt{sample\_count} & 200 & 200.0 & 200 \\
\hline
\texttt{sampling\_frequency\_hz} & 40.0 & 40.0 & 40.0 \\
\hline
\texttt{dominant\_frequency\_hz} & 5.0 & 5.0 & 5.0 \\
\hline
\texttt{payload\_size\_bytes} & 454 & 454.0 & 454 \\
\hline
\end{tabular}
\end{center}

\subsection{Interpretation}

The edge path is no longer just implemented in code. It has been validated
end-to-end with:

\begin{itemize}
    \item real board,
    \item real WiFi,
    \item real broker,
    \item real listener output,
    \item saved artifacts.
\end{itemize}

\reportfigure{\textwidth}{../results/screenshots/mqtt_listener_received_message_2026-04-18.png}{Python edge listener receiving consecutive MQTT aggregate messages from the Heltec board.}

\section{Three-Signal Validation}

The bonus-ready signal profiles were also validated through real-board MQTT
runs.

\subsection{Comparison Table}

\begin{center}
\small
\resizebox{\textwidth}{!}{
\begin{tabular}{|l|r|r|r|r|r|}
\hline
\textbf{Profile} & \textbf{Avg dominant [Hz]} & \textbf{Avg sampling [Hz]} & \textbf{Avg anomaly} & \textbf{Avg payload [B]} & \textbf{Avg listener latency [us]} \\
\hline
\texttt{clean\_reference} & 5.000 Hz & 40.000 Hz & 0.000 & 446.0 B & 854014.8 us \\
\hline
\texttt{noisy\_reference} & 5.000 Hz & 40.000 Hz & 0.000 & 445.6 B & 845563.8 us \\
\hline
\texttt{anomaly\_stress} & 5.000 Hz & 40.000 Hz & 2.800 & 444.6 B & 823576.4 us \\
\hline
\end{tabular}
}
\normalsize
\end{center}

\subsection{Main Observations}

\begin{itemize}
    \item all three profiles preserved \texttt{5.00 Hz} as the dominant frequency in the saved windows,
    \item the anomaly profile produced non-zero \texttt{anomaly\_count} values as expected,
    \item the payload size remained almost constant because each window still produces one compact aggregate.
\end{itemize}

\reportfigure{\textwidth}{../results/screenshots/signal_profile_comparison_2026-04-18.png}{Comparison of the three saved signal-profile runs captured from the real board over MQTT.}

\reportfigure{0.9\textwidth}{../results/screenshots/anomaly_profile_nonzero_count_2026-04-17.png}{Non-zero anomaly counts recorded in the anomaly-stress profile.}

\section{BetterSerialPlotter Validation}

The project includes a simplified serial stream for \texttt{BetterSerialPlotter}.
This visual evidence shows:

\begin{itemize}
    \item current sampling frequency,
    \item dominant frequency,
    \item aggregate average.
\end{itemize}

This figure is useful for presenting the live signal-processing pipeline, but
it is not used as proof of MQTT or TTN delivery.

\reportfigure{\textwidth}{../pics/2026-04-18_better_serial_plotter_live_view.png}{BetterSerialPlotter receiving the live adaptive-sampling stream from the Heltec board.}

\section{LoRaWAN And TTN Status}

The \texttt{LoRaWAN} path is implemented up to the payload-preparation stage.

What is already true:

\begin{itemize}
    \item the aggregate is serialized into a compact uplink payload,
    \item the LoRa path is integrated into the same pipeline as MQTT,
    \item the runtime logs show prepared payloads such as \texttt{000100C8019001F4FFE8}.
\end{itemize}

What is still missing:

\begin{itemize}
    \item a real uplink through a reachable gateway,
    \item \texttt{TTN} live data screenshots,
    \item decoded cloud-side proof.
\end{itemize}

So the correct current description is:

\begin{itemize}
    \item \texttt{LoRaWAN/TTN} path is implemented and stub-ready,
    \item live \texttt{TTN} validation is still pending.
\end{itemize}

\section{Secure MQTT Status}

The MQTT client now supports:

\begin{itemize}
    \item plain \texttt{mqtt://},
    \item verified \texttt{mqtts://},
    \item optional username/password authentication,
    \item certificate verification through the ESP-IDF CA bundle.
\end{itemize}

However, the captured saved run is still the plain local-broker run.

So the honest current status is:

\begin{itemize}
    \item secure MQTT support is implemented in firmware,
    \item a live secure-broker validation run is still pending.
\end{itemize}

\section{Energy Measurement Status}

Energy is a required grading item, but the cleanest measurement setup needs two
boards:

\begin{itemize}
    \item one board under test,
    \item one board or monitor path reading the \texttt{INA219}.
\end{itemize}

Because only one Heltec is currently available, the project keeps the energy
section prepared but deferred.

What has already been done:

\begin{itemize}
    \item the firmware now has a baseline-versus-adaptive switch,
    \item the measurement runbook has been written,
    \item the \texttt{INA219} setup for the Heltec V3 has been documented.
\end{itemize}

What is still missing:

\begin{itemize}
    \item real measured power or energy values,
    \item screenshots of the measurement setup,
    \item final baseline-versus-adaptive comparison table.
\end{itemize}

\section{Communication Volume Discussion}

One subtle but important observation is that this project sends one aggregate
per completed window, not raw samples over MQTT.

That means:

\begin{itemize}
    \item adaptive sampling reduces local sampling and processing work,
    \item but it does not automatically reduce the MQTT message size very much.
\end{itemize}

This is exactly what the saved payload summaries show:

\begin{itemize}
    \item payload sizes remain around \texttt{444} to \texttt{454} bytes depending on the profile.
\end{itemize}

So the communication-volume section in the final report should explain that:

\begin{itemize}
    \item the main gain from adaptive sampling in this implementation is local efficiency,
    \item the network payload count per window is mostly unchanged.
\end{itemize}

\section{Build And Resource Status}

The current firmware still builds successfully after the recent additions such
as:

\begin{itemize}
    \item three signal profiles,
    \item secure MQTT support,
    \item baseline/adaptive configuration switch.
\end{itemize}

Current build snapshot:

\begin{center}
\begin{tabular}{|l|l|}
\hline
\textbf{Resource} & \textbf{Usage} \\
\hline
RAM & \texttt{38520 B} (\texttt{11.8\%}) \\
\hline
Flash & \texttt{973469 B} (\texttt{92.8\%}) \\
\hline
\end{tabular}
\end{center}

This is still acceptable, but flash headroom is now limited.

\section{Remaining Work}

The remaining work before final submission is:

\begin{center}
\begin{tabularx}{\textwidth}{|>{\raggedright\arraybackslash}p{0.23\textwidth}|>{\raggedright\arraybackslash}p{0.14\textwidth}|X|}
\hline
\textbf{Item} & \textbf{Status} & \textbf{Notes} \\
\hline
\texttt{LoRaWAN/TTN} live proof & pending & must be done on campus or under real gateway coverage \\
\hline
Energy measurements & pending & deferred until a cleaner two-board setup is available \\
\hline
Secure MQTT live TLS proof & pending & firmware support exists, but no saved live TLS run yet \\
\hline
Prompt log curation & pending & intentionally left for later \\
\hline
Final short report & pending & this detailed report will make that much easier \\
\hline
\end{tabularx}
\end{center}

\section{Conclusion}

At this checkpoint, the project has moved well beyond the early scaffolding
phases.

The most important completed achievements are:

\begin{itemize}
    \item stable virtual-signal generation on the real board,
    \item measured maximum stable sampling benchmark,
    \item correct dominant-frequency detection,
    \item working adaptive-sampling transition from \texttt{50 Hz} to \texttt{40 Hz},
    \item correct per-window aggregate computation,
    \item full home-network \texttt{MQTT/WiFi} validation on hardware,
    \item real latency measurements based on synchronized timestamps,
    \item three validated signal profiles including anomaly-aware evidence.
\end{itemize}

The most important missing items are now concentrated and clear:

\begin{itemize}
    \item real \texttt{TTN} uplink evidence,
    \item energy comparison measurements,
    \item one secure \texttt{MQTTS} validation run.
\end{itemize}

This means the project is already in a strong state for the workshop: the core
pipeline works, the edge path is proven, and the remaining work is focused on
final measurement and cloud validation rather than on missing core
functionality.

\end{document}
